\begin{definition}[Fixed two-plaquette graph and tree]
    \label{def:peps_two_plaquette_graph}
    \lean{TNLean.PEPS.twoPlaquetteGraph,
        TNLean.PEPS.twoPlaquetteTree,
        TNLean.PEPS.twoPlaquetteCycleBond}
    \leanok
    Number the six sites in the order
    $(0,0),(1,0),(2,0),(2,1),(1,1),(0,1)$.
    Let $T$ be the path joining consecutive numbered sites, and let
    $\Gamma$ add to $T$ the bonds $e_0=\{0,5\}$ and $e_1=\{1,4\}$.
    These are respectively the leftmost and middle vertical bonds.
    This is the six-site geometry used in the construction of
    \cite[proof of moving fluxons]{Schuch2010PEPS}.
\end{definition}

\begin{theorem}[The five-edge tree and two remaining bonds]
    \label{thm:peps_two_plaquette_graph}
    \lean{TNLean.PEPS.twoPlaquetteTree_le,
        TNLean.PEPS.twoPlaquetteTree_isTree,
        TNLean.PEPS.twoPlaquetteGraph_connected,
        TNLean.PEPS.twoPlaquetteGraph_card_edges,
        TNLean.PEPS.twoPlaquetteTree_card_edges,
        TNLean.PEPS.twoPlaquetteCycleBond_not_tree,
        TNLean.PEPS.twoPlaquetteCycleBond_injective,
        TNLean.PEPS.twoPlaquetteCycleBond_surjective_non_tree}
    \leanok
    \uses{def:peps_two_plaquette_graph}
    The graph $\Gamma$ is connected and has seven bonds.
    The subgraph $T$ is a spanning tree with five bonds.
    The bonds $e_0,e_1$ are distinct, and they are exactly the
    bonds of $\Gamma$ outside $T$.
\end{theorem}
\begin{proof}\leanok
    The path has six vertices and five consecutive bonds, hence is a tree.
    Adding the two stated chords gives the two-plaquette graph.
    The endpoint pairs enumerate all seven bonds; exactly the two chords
    lie outside the path.
\end{proof}

\begin{definition}[Placement of the two-plaquette block]
    \label{def:peps_two_plaquette_torus_geometry}
    \lean{TNLean.PEPS.twoPlaquetteTorusVertex,
        TNLean.PEPS.twoPlaquetteTorusRegion,
        TNLean.PEPS.twoPlaquetteTorusVertexEquiv,
        TNLean.PEPS.twoPlaquetteTorusIso,
        TNLean.PEPS.twoPlaquetteTorusTree,
        TNLean.PEPS.twoPlaquetteTorusCycleBond,
        TNLean.PEPS.twoPlaquetteTorusVertex_injective,
        TNLean.PEPS.twoPlaquetteTorusVertex_adj_iff}
    \leanok
    \uses{def:peps_two_plaquette_graph,
        thm:peps_two_plaquette_graph,thm:peps_edge_map_subgraph}
    On a torus of horizontal period at least four and vertical period
    at least three, place the six sites at the displayed coordinates.
    Their numbering identifies $\Gamma$ with the actual induced torus
    graph on the resulting region $R$. Transport $T,e_0,e_1$ by this
    identification, retaining the ordered bond endpoints.
\end{definition}

\begin{theorem}[Actual six-site torus geometry]
    \label{thm:peps_two_plaquette_torus_geometry}
    \lean{TNLean.PEPS.twoPlaquetteTorusTree_isTree,
        TNLean.PEPS.twoPlaquetteTorusTree_le,
        TNLean.PEPS.twoPlaquetteTorusRegion_card,
        TNLean.PEPS.twoPlaquetteTorusRegion_eq_rectangle,
        TNLean.PEPS.twoPlaquetteTorusCycleBond_injective,
        TNLean.PEPS.twoPlaquetteTorusRegion_card_internalEdges,
        TNLean.PEPS.twoPlaquetteTorusTree_card_edges,
        TNLean.PEPS.twoPlaquetteTorusRegion_connected,
        TNLean.PEPS.twoPlaquetteTorusCycleBond_endpoints}
    \leanok
    \uses{def:peps_two_plaquette_torus_geometry,
        thm:peps_two_plaquette_graph}
    The actual region $R$ is the coordinate rectangle
    $\{0,1,2\}\times\{0,1\}$. It has six physical sites,
    seven internal bonds, and a connected induced graph.
    The transported $T$ is a spanning tree with five bonds.
    Its two selected non-tree bonds are distinct and have ordered endpoints
    $(0,0)\to(0,1)$ and $(1,0)\to(1,1)$.
    These statements concern the ambient torus region, including its exterior
    boundary.
\end{theorem}
\begin{proof}\leanok
    All horizontal coordinates lie between zero and two; all vertical
    coordinates lie between zero and one. The period bounds exclude
    additional wraparound adjacency inside the region. Thus coordinate
    comparison preserves and reflects each bond of $\Gamma$.
    The resulting graph isomorphism transports the spanning tree,
    connectivity, and both edge counts. The lexicographic order
    orients each selected vertical bond upward.
\end{proof}
